%----------------------------------------------------------------------------------------
%	PACKAGES AND OTHER DOCUMENT CONFIGURATIONS
%----------------------------------------------------------------------------------------



\documentclass[a4paper,12pt]{article}
\usepackage[english]{babel}
\usepackage[latin1]{inputenc}
\usepackage{amsmath}
\usepackage{amssymb}
\usepackage{amsfonts}
\usepackage{graphicx}
\usepackage[colorinlistoftodos]{todonotes}
\usepackage[toc,page]{appendix}
\usepackage{setspace}
\doublespacing
\usepackage{booktabs}
\usepackage{geometry}
\usepackage[bottom]{footmisc}
\usepackage{longtable}
\usepackage[demo]{graphicx}
\usepackage{subfig}
\usepackage{multirow}
\renewcommand{\arraystretch}{0.7}
\renewcommand{\labelitemi}{$\triangleright$}
 \geometry{
 a4paper,
 total={170mm,257mm},
 left=25mm,
 top=30mm,
 right=25mm,
 bottom=25mm,
 }
 \usepackage{hyperref}
 \hypersetup{
    bookmarks=true,         % show bookmarks bar?
    unicode=false,          % non-Latin characters in Acrobat’s bookmarks
    pdftoolbar=true,        % show Acrobat’s toolbar?
    pdfmenubar=true,        % show Acrobat’s menu?
    pdffitwindow=false,     % window fit to page when opened
    pdfstartview={FitH},    % fits the width of the page to the window
    pdftitle={Design_Report},    % title
    pdfauthor={Jacob Pichelman, Luca Poll},     % author
    pdfsubject={Subject},   % subject of the document
    pdfcreator={Creator},   % creator of the document
    pdfproducer={Producer}, % producer of the document
    pdfkeywords={keyword1, key2, key3}, % list of keywords
    pdfnewwindow=true,      % links in new PDF window
    colorlinks=false,       % false: boxed links; true: colored links
    linkcolor=red,          % color of internal links (change box color with linkbordercolor)
    citecolor=green,        % color of links to bibliography
    filecolor=magenta,      % color of file links
    urlcolor=cyan           % color of external links
}

	% biblatex
\usepackage[
       style=authoryear,
       natbib=true,
       maxcitenames=3,
       maxbibnames=11,
       backend=biber,
       pagetracker=page,
       hyperref=true,
       doi=true, 
       mergedate=compact, 
       firstinits=true
    ]{biblatex} 
	\usepackage{csquotes}
	\renewcommand*{\bibsetup}{
		\interlinepenalty=10000\relax % default is 5000
		\widowpenalty=10000\relax
		\clubpenalty=10000\relax
		\raggedbottom
		\frenchspacing
        \biburlsetup}
    \addbibresource{library.bib}

\begin{document}
\begin{titlepage}

\newcommand{\HRule}{\rule{\linewidth}{0.25mm}} % Defines a new command for the horizontal lines, change thickness here
\setlength{\topmargin}{-0.5in}
\center % Center everything on the page


\includegraphics[scale=0.75]{TSE.png}\\



%----------------------------------------------------------------------------------------
%	HEADING SECTIONS
%----------------------------------------------------------------------------------------
\\[1.5cm]
\large \textsc{M2 EEE Panel Data} 
\vspace{1.5cm}
% Name of your heading such as course name
\textsc{\large } % Minor heading such as course title


%----------------------------------------------------------------------------------------
%	TITLE SECTION
%----------------------------------------------------------------------------------------

\HRule \\[0.75cm]
{ \huge \bfseries Panel Data Research Proposal}\\[0.5cm] % Title of your document
\HRule \\[1.75cm]
 
%----------------------------------------------------------------------------------------
%	AUTHOR SECTION
%----------------------------------------------------------------------------------------


\large\textsc{Andrew Boomer, \\ Jacob Pichelmann, \\Luca Poll} \\[1.5cm]

%----------------------------------------------------------------------------------------
%	DATE SECTION
%----------------------------------------------------------------------------------------

{\large \today}\\[0.5cm] % Date, change the \today to a set date if you want to be precise

\vfill % Fill the rest of the page with whitespace

\end{titlepage}


%-------------------------------------------------------------

% TABLE OF CONTENT
\renewcommand{\contentsname}{Table of Contents}
\tableofcontents
\clearpage
%-----------------------------------------------------------------------------



\section{Introduction}
After the Arab spring and the related outbreak of unforeseen violence, conflict forecasting models were largely criticized and it was argued, that forecasting new civil wars might have reached a limit. Mueller and Rauh (2018) though show in their paper "Reading between the lines: Prediction of political violence", that this might not be entirely true. Their main argument is structured as follows: Conventional conflict forecasting models\footnote{\noindent They demonstrate their argument by replicating the following papers on conflict prediction: \begin{itemize}
    \item Miguel \& Satyanath (2011): Prediction through rainfall growth
    \item Besley \& Presson (2011): Prediction through proxies for external shocks and political constraints
    \item Goldstone et al. (2010): Prediction through political institution dummies, child mortality rates, share of population discriminated against and whether neighboring countries in conflict
    \item Ward et al. (2013): Event database on high-intensity and low-intensity conflict events used for analysis
    \item Chadefaux (2014): Conflict prediction through analysis of keyword count in newspaper text
\end{itemize}}, that rely on the overall variation in country fixed effect models, exhibit a bias towards predicting conflict onset to where conflict has occurred before. This is partially due to large country fixed effects and slow moving factors like population, ethnic fractionalization, climate, etc. that result in a large between variation. The forecasts are hence dominated by structural time-invariant (or slow moving) factors, neglecting valuable within variation. As a result these models are relatively good at predicting (biasedly) where conflict will happen, but not when it will happen. In order to improve the forecasting of the timing of conflict and generate an unbiased forecast, Mueller \& Rauh (2018) propose to isolate the within from the overall variation and use such to predict the onset of armed conflict and civil war. In order to obtain necessary within variation, they propose to use topic modeling on newspaper text and create variables informing about the average distribution of topic shares observed in a country during one year. 


\section{Sample \& Data}

The sample for the underlying empirical analysis consists of 700.000 newspaper articles from three internationally-reporting newspapers between 1975 and 2015: the Economist\footnote{174.450 articles from 1975 onward}, the New York Times\footnote{363.275 articles from 1980 onward} and the Washington Post\footnote{185.523 articles from 1977 onward}.
The newspapers cover in total 185 countries and the average yearly coverage amounts to 120 articles per country (with a range from 1 to 5.500). \\
The dependent variables on the other hand are constructed through battle-related deaths from the Uppsala Conflict Data Program (UCDP/PRIO). Following their definition, armed conflict (dep. var. 1) is defined as a contested incompatibility that concerns government and/or territory over which the use of armed force between two parties, of which at least one is the government of a state, has resulted in at least 25 battle-related deaths in one calendar year. Civil conflict (dep. var. 2) follows the same definition but requires at least 1.000 battle-related deaths in on calendar year. 




\section{Model}
The aim of the model is to create forecasts for an armed conflict/ civil war outbreak in period $T+1$ at period $T \in \{1995,..., 2013\}$. To create this forecast, every available information up to period $T$ is included into the forecast. Therefore, the respective country-year topic shares $\theta_{n,i,T}$ are calculated for every newspaper sub-sample available up to period $T$\footnote{As the amount of available articles/ words expands in $T$, the basis for defining a topic through characteristic words in $T$ does also expand. Hence, the every topic characteristic and every topic distribution will vary at every $T$} for each country $i$ and topic $n$. As a consequence, the following two steps are repeated at every $T$:

\noindent\textbf{Step 1: Estimate model and obtain fitted values}

\noindent From the model $y_{i,T+1} = \alpha + \beta_{i} + \theta_{i,T}\beta^{topics}$ the fitted values from the estimation based on the overall variation are obtained: 

\begin{equation}
    \hat{y}_{i,T+1}^{overall} = \hat{\alpha} + \hat{\beta_i} + \theta_{i,T}\hat{\beta}^{topics}
\end{equation}

\noindent From these fitted values that rely on the overall variation, the fitted fixed effects are subtracted in order to obtain the fitted within model:

\begin{equation}
    \hat{y}_{i,T+1}^{within} = \hat{\alpha} + \theta_{i,T}\hat{\beta}^{topics}
\end{equation}

\noindent \textbf{Step 2: Produce forecast based on fitted values for period T+1}

\noindent 1) The fitted values are transformed into binary variables depending on cutoff value c\\
2) Compare forecast (binary variable) to realizations of armed conflict and civil war\\
3) Assess performance of overall and within model by considering forecasting performance for any given value c through ROC curves 


\section{Extensions}
After replicating the core findings of the paper, we plan to extend it in several dimensions, following a set of four steps. 

\subsection{Step I: Extended Data Analysis}
Firstly, we will place emphasis on investigating the characteristics of the data at hand, i.e. analyzing the distribution of topic shares over time and their correlation with certain confounders. We control for possible correlations with country specific characteristics such as continents/regions and rainfall and check whether or not there exist outliers (e.g. countries with just one article a year or where coverage only occurred during conflict). We aim to build a sound argument for the resulting estimation method and try to answer if additional measures, such as employing panel robust standard errors, %Note: I was just thinking that they actually dont really care (or at least never talk about it) about the SE or significance for creating the binary variable.... Maybe we can do sth there too
should be taken into consideration. We expect to additionally validate the authors' econometric approach with this exercise. 

\subsection{Step II: Enhanced Estimation}
Based on the findings in step I we advance by employing different models in order to increase the estimates precision. Depending on said findings this will include: 
\begin{itemize}
    \item Conditional logit - including a discussion of necessary (strict) assumptions (i.e. fixed effects of a country that either always or never experience conflict)
    \item Random Effects - the assumption that $\alpha_i$ is uncorrelated with the regressors (i.e. topic shares) is unlikely to hold. However, when interpreting the random effects estimator as a weighted average of within and between estimation its performance in conjunction with the results of the fixed effects estimation might allow for a clearer assessment of the underlying factors at play. Additionally, the random effects model allows for a straightforward out of sample prediction.
    \item First Differences - relaxing the strict exogeneity assumption of the fixed effects model
\end{itemize}

\subsection{Step III: Data Changes}
The key goal of the paper is to develop a forecasting strategy that allows to predict conflict in formerly peaceful regions, i.e. a regime change. We plan to additionally assess our model performance by building another dependent variable which captures regime change itself and regress on this newly constructed variable.

% Moreover, there might be instances where a more timely prediction of conflict is of use. We will hence change the time horizon from yearly to quarterly. %%%% NOT POSSIBLE DUE TO SIZE OF SAMPLE %%%%%

Moreover, we plan to include further lags of topics in order to capture a possible 'build up' of conflict over the preceding years. However, the topic composition might lack the level of granularity necessary to capture these patterns. 

\subsection{Step IV: Enhanced discussion}
Finally our analysis aims at discovering possible additional sources of bias, to better understand model performance. However, given the data constraints, we do not expect to resolve these issues but rather add to the discussion. 
For example, we implicitly assume that the newspapers have unbiased reporting. This assumption is likely not to hold. We will investigate options to control for an author specific bias. 
Additionally, we will further investigate the (changing) importance of topics over the years. This exercise aims at distilling a set of topics that is predictive for each time span.



\section*{References}
Besley, Timothy and Torsten Persson. 2011. Pillars of prosperity: The political economics \newline \indent of development clusters. Princeton University Press.\vspace{0.25cm}

\noindent Chadefaux, Thomas. 2014. "Early warning signals for war in the news." Journal of Peace \newline \indent Research 51(1):5-18.\vspace{0.25cm}

\noindent Goldstone, Jack A, Robert H Bates, David L Epstein, Ted Robert Gurr, Michael B Lustik, \newline \indent Monty G Marshall, Jay Ulfelder and Mark Woodward. 2010. "A global model for \newline \indent forecasting political instability." American Journal of Political Science 54(1):190-208.\vspace{0.25cm}

\noindent Miguel, Edward and Shanker Satyanath. 2011. "Re-examining economic shocks and civil \newline \indent conflict." American Economic Journal: Applied Economics 3(4):228-232.\vspace{0.25cm}

\noindent Mueller, H., & Rauh, C. (2018). "Reading Between the Lines: Prediction of Political \newline \indent Violence Using Newspaper Text." American Political Science Review, 112(2), 358-375. \newline \indent doi:10.1017/S0003055417000570\vspace{0.25cm}

\noindent Ward, Michael D, Nils W Metternich, Cassy L Dorff, Max Gallop, Florian M Hollenbach, \newline \indent Anna Schultz and Simon Weschle. 2013. "Learning from the past and stepping into \newline \indent the future: Toward a new generation of conflict prediction." International Studies \newline \indent Review 15(4):473-490.

\clearpage

\subsection*{Summary Statistics of Panel for Full Sample (2013)}
\begin{longtable}{lllrrrr}

\toprule
               &        &  Mean &  Std. Dev. &    Min &   Max &  Observations \\
\textbf{Variable} & \textbf{Type} &       &            &        &       &               \\
\midrule
\multirow{3}{*}{\textbf{Armed Conflict}} & \textbf{overall} & 0.142 &      0.349 &  0.000 & 1.000 &          7520 \\
               & \textbf{between} &       &      0.020 &  0.106 & 0.186 &            40 \\
               & \textbf{within} &       &      0.349 & -0.044 & 1.036 &           188 \\
\cline{1-7}
\multirow{3}{*}{\textbf{Civil War}} & \textbf{overall} & 0.060 &      0.237 &  0.000 & 1.000 &          7520 \\
               & \textbf{between} &       &      0.024 &  0.027 & 0.112 &            40 \\
               & \textbf{within} &       &      0.236 & -0.052 & 1.033 &           188 \\
\cline{1-7}
\multirow{3}{*}{\textbf{Topic 1 Share}} & \textbf{overall} & 0.053 &      0.039 &  0.007 & 0.560 &          6639 \\
               & \textbf{between} &       &      0.005 &  0.046 & 0.063 &            39 \\
               & \textbf{within} &       &      0.038 & -0.002 & 0.561 &           185 \\
\cline{1-7}
\multirow{3}{*}{\textbf{Topic 2 Share}} & \textbf{overall} & 0.073 &      0.041 &  0.010 & 0.559 &          6639 \\
               & \textbf{between} &       &      0.010 &  0.050 & 0.089 &            39 \\
               & \textbf{within} &       &      0.040 &  0.004 & 0.549 &           185 \\
\cline{1-7}
\multirow{3}{*}{\textbf{Topic 3 Share}} & \textbf{overall} & 0.043 &      0.049 &  0.006 & 0.454 &          6639 \\
               & \textbf{between} &       &      0.003 &  0.038 & 0.051 &            39 \\
               & \textbf{within} &       &      0.049 &  0.004 & 0.451 &           185 \\
\cline{1-7}
\multirow{3}{*}{\textbf{Topic 4 Share}} & \textbf{overall} & 0.060 &      0.068 &  0.009 & 0.663 &          6639 \\
               & \textbf{between} &       &      0.012 &  0.032 & 0.080 &            39 \\
               & \textbf{within} &       &      0.067 & -0.006 & 0.663 &           185 \\
\cline{1-7}
\multirow{3}{*}{\textbf{Topic 5 Share}} & \textbf{overall} & 0.069 &      0.045 &  0.004 & 0.468 &          6639 \\
               & \textbf{between} &       &      0.008 &  0.045 & 0.081 &            39 \\
               & \textbf{within} &       &      0.045 & -0.003 & 0.476 &           185 \\
\cline{1-7}
\multirow{3}{*}{\textbf{Topic 6 Share}} & \textbf{overall} & 0.063 &      0.052 &  0.009 & 0.765 &          6639 \\
               & \textbf{between} &       &      0.011 &  0.036 & 0.081 &            39 \\
               & \textbf{within} &       &      0.051 & -0.003 & 0.774 &           185 \\
\cline{1-7}
\multirow{3}{*}{\textbf{Topic 7 Share}} & \textbf{overall} & 0.074 &      0.047 &  0.007 & 0.514 &          6639 \\
               & \textbf{between} &       &      0.006 &  0.063 & 0.086 &            39 \\
               & \textbf{within} &       &      0.046 & -0.005 & 0.509 &           185 \\
\cline{1-7}
\multirow{3}{*}{\textbf{Topic 8 Share}} & \textbf{overall} & 0.070 &      0.052 &  0.007 & 0.426 &          6639 \\
               & \textbf{between} &       &      0.006 &  0.058 & 0.084 &            39 \\
               & \textbf{within} &       &      0.051 & -0.006 & 0.420 &           185 \\
\cline{1-7}
\multirow{3}{*}{\textbf{Topic 9 Share}} & \textbf{overall} & 0.074 &      0.054 &  0.010 & 0.514 &          6639 \\
               & \textbf{between} &       &      0.012 &  0.058 & 0.116 &            39 \\
               & \textbf{within} &       &      0.053 & -0.024 & 0.519 &           185 \\
\cline{1-7}
\multirow{3}{*}{\textbf{Topic 10 Share}} & \textbf{overall} & 0.065 &      0.051 &  0.007 & 0.612 &          6639 \\
               & \textbf{between} &       &      0.008 &  0.053 & 0.092 &            39 \\
               & \textbf{within} &       &      0.051 & -0.009 & 0.605 &           185 \\
\cline{1-7}
\multirow{3}{*}{\textbf{Topic 11 Share}} & \textbf{overall} & 0.063 &      0.046 &  0.005 & 0.407 &          6639 \\
               & \textbf{between} &       &      0.010 &  0.047 & 0.082 &            39 \\
               & \textbf{within} &       &      0.044 & -0.008 & 0.410 &           185 \\
\cline{1-7}
\multirow{3}{*}{\textbf{Topic 12 Share}} & \textbf{overall} & 0.075 &      0.069 &  0.004 & 0.653 &          6639 \\
               & \textbf{between} &       &      0.017 &  0.058 & 0.135 &            39 \\
               & \textbf{within} &       &      0.067 & -0.044 & 0.654 &           185 \\
\cline{1-7}
\multirow{3}{*}{\textbf{Topic 13 Share}} & \textbf{overall} & 0.089 &      0.090 &  0.008 & 0.623 &          6639 \\
               & \textbf{between} &       &      0.010 &  0.070 & 0.103 &            39 \\
               & \textbf{within} &       &      0.090 & -0.001 & 0.614 &           185 \\
\cline{1-7}
\multirow{3}{*}{\textbf{Topic 14 Share}} & \textbf{overall} & 0.067 &      0.048 &  0.007 & 0.582 &          6639 \\
               & \textbf{between} &       &      0.005 &  0.058 & 0.076 &            39 \\
               & \textbf{within} &       &      0.048 &  0.006 & 0.579 &           185 \\
\cline{1-7}
\multirow{3}{*}{\textbf{Topic 15 Share}} & \textbf{overall} & 0.061 &      0.055 &  0.006 & 0.437 &          6639 \\
               & \textbf{between} &       &      0.007 &  0.048 & 0.075 &            39 \\
               & \textbf{within} &       &      0.055 & -0.006 & 0.429 &           185 \\
\bottomrule
\caption{Panel Data Summary}
\label{tab:xtsum}
\end{longtable}


\end{document}